
\chapter{Conclusion}
\label{chap:conclusion}

\section{Summary of Findings}
\label{sec:concl-summary}

This thesis investigated how a frozen latent world model can be used to
improve offline reinforcement learning for high-dimensional robotic
manipulation. We evaluated a latent action-conditioned world model in several roles: as
a training environment, as a short-horizon planner, as a source of
world-model-generated training data, and as a tool for grounding
hierarchical subgoals. The central finding is that these roles are not
equivalent. The world model is useful when used locally and
conservatively, but unreliable when treated as a replacement for the
environment or as an unconstrained source of training data.

Using the world model as a surrogate environment for reinforcement
learning consistently failed to improve over the offline baseline.
Across sparse rewards, dense rewards, anchoring, joint
predictor--policy training, uncertainty penalties, and
nearest-neighbour regularisation, policy optimisation inside the
learned dynamics either underperformed or collapsed. The failure mode is
consistent with the predictor-drift analysis: small latent prediction
errors accumulate over rollout horizons, and once the critic is allowed
to bootstrap on imagined successors, value estimates become corrupted.
This shows that accurate short-horizon prediction is not sufficient for
stable long-horizon reinforcement learning inside a learned simulator.

In contrast, the world model is useful for short-horizon
value-guided planning. Best-of-$N$ MPC improves the single-task offline
policy from $81.6\%$ to $88.0\%$, and gradient-based value-guided
planning reaches $90.0\%$ with no parameter updates. FMQ achieves the
same level on the frozen offline critic, but only within a narrow
trust-region range. After training the actor and critic on
planner-generated actions, FMQ inference becomes substantially more
robust and reaches $96.0\%$ success. This result highlights the key
mechanism: the planner benefits from the world model's local
counterfactual predictions, but only when the critic remains calibrated
to the actions being evaluated.

The online-in-WM experiments provide the clearest evidence for how
imagined transitions should be used during training. Updating both actor
and critic on world-model-generated transitions collapses to
$30.7 \pm 20.8\%$. Freezing the critic and updating only the actor
instead reaches a standalone peak of $94.8 \pm 0.7\%$ and a final
success rate of $90.7 \pm 0.8\%$ on the single-task benchmark. The
world model can therefore provide useful policy-improvement
supervision, but should not be used to generate bootstrapping targets
for value learning.

The same recipe does not transfer cleanly to the full
goal-conditioned benchmark. Goal-conditioned online-in-WM briefly
matches or exceeds the offline baseline, but final performance remains
below it; the strongest leakage-free achieved-state-goal variant reaches
$84.4 \pm 0.9\%$ compared with a $90.0\%$ offline baseline. Generating a
fixed offline corpus from the world model also underperforms real play
data: with identical HER processing, FQL reaches $90.3 \pm 1.6\%$ on
real data but only $64.7 \pm 3.7\%$ on the best world-model-generated
corpus. These results suggest that the successful single-task methods
depend on closed-loop interaction between the current actor, planner,
and frozen critic. Once that loop is broken, much of the benefit is
lost.

Latent Subgoal Planning (LSP) provides the strongest stable
goal-conditioned result. Phase 1, trained entirely offline with
flow-matching high- and low-level policies, reaches $88.0\%$ on the
five-task benchmark, slightly exceeding the HIQL end-to-end baseline.
However, Phase 2 attempts to improve the high-level policy using
world-model-grounded AWR or achieved-state subgoal regression do not
surpass this baseline. Conservative updates preserve some Phase-1
performance but do not improve it, while more aggressive updates
collapse the hierarchy. The likely cause is distributional coupling:
shifting the high-level subgoal distribution pushes the frozen low-level
policy outside the conditioning distribution it was trained to execute.

Several early SAC-based approaches, including curriculum SAC, HER+SAC,
and hierarchical SAC with expert-guided stabilisation, were also
explored. These methods could often optimise latent-reaching objectives
inside the world model, but did not transfer to native task success in
MuJoCo. They are therefore treated as exploratory negative results and
documented separately in the appendix.

Overall, the thesis supports a constrained view of learned world models
for offline robotic control. The world model is not reliable enough to
serve as a standalone simulator, a replacement offline dataset, or an
unconstrained subgoal generator. Its strength lies in short-horizon,
local counterfactual reasoning. When paired with a critic trained on
real transitions, this local reasoning can improve action selection,
generate useful actor-training data, and produce substantial gains in
single-task manipulation.

\section{Future Work}
\label{sec:concl-future}

\paragraph{Improving world-model reliability and action conditioning.}
The methods that succeed in this thesis all avoid long-horizon open-loop
rollouts. They use the world model over one or two prediction steps and
rely on a critic trained on real data to score the result. Future work
should repeat the same diagnostic pipeline with world-model
architectures that may offer longer reliable horizons or stronger action
conditioning, such as diffusion-transformer world models,
Dreamer-style recurrent state-space models, or value-aware latent
models. The action-conditioning ratio used in the generalisation
experiments is also a useful pre-screening diagnostic: online-in-WM
should only be expected to help when the learned dynamics are
sufficiently sensitive to actions.

\paragraph{Scaling generated-data methods to goal-conditioned control.}
The single-task online-in-WM method improves substantially, but the
goal-conditioned extension does not reliably beat the offline baseline.
Future work should focus on coverage: more diverse goal sampling,
explicit exploration bonuses in the MPC scorer, higher weighting of
real offline replay during online updates, or goal-conditioned critics
specialised to different regions of the goal space. These directions
target the main weakness observed here: the generated online buffer
becomes concentrated around the states visited by the current planner,
rather than covering the full state-goal distribution.

\paragraph{Stabilising hierarchical world-model learning.}
The LSP Phase-2 failures leave a clear open problem: how can a
hierarchical policy use world-model-derived supervision without breaking
coordination between the high and low levels? In the flat single-task
case, the successful rule is to freeze the critic and update only the
actor. The analogous rule for hierarchical agents is less obvious,
because changing the high-level policy also changes the conditioning
distribution faced by the low-level policy. Future work could explore
conservative subgoal trust regions, explicit low-level reachability
models, or joint update rules that maintain compatibility between
proposed subgoals and low-level competence. The broader challenge is to extend the local usefulness of learned
world models to broader goal-conditioned and hierarchical settings
without losing the grounding provided by real experience.

